\subsection*{One group, one tile: the simplest MMA matmul}

The smallest example assigns one output tile to a cooperative group, written \co@: group@. It accepts left and right operands plus an output buffer supplied by the caller. Follow one output tile through the code: initialize its accumulator, visit every K slice, store the accumulated tile in shared storage, and copy it to its output position.

\begin{lstlisting}[language=croqtile,caption={One group, one tile: the simplest MMA matmul}]
#define TILE_M 16
#define TILE_N 16
#define TILE_K 16
#define WARP_M 16
#define WARP_N 16
#define WARP_K 16
__co__ void matmul_wmma(global f16 [M, K] lhs, global f16 [N, K] rhs, global f16 [M, N] output) {
  parallel {block_m, block_n} by [cdiv(M, TILE_M), cdiv(N, TILE_N)] : block {
    shared f16 [TILE_M, TILE_N] output_s;
    parallel {warp_m, warp_n} by [cdiv(TILE_M, MMA_M), cdiv(TILE_N, MMA_N)] : group {
      mc = mma.fill 0.0;
      foreach {iv_k} in [cdiv(K, TILE_K)] {
        lhs_load_s = dma.copy lhs.subspan(TILE_M, TILE_K).at(block_m, iv_k) => shared;
        rhs_load_s = dma.copy rhs.subspan(TILE_N, TILE_K).at(block_n, iv_k) => shared;
        foreach iv_warp_k in [cdiv(TILE_K, MMA_K)] {
          ma = mma.load lhs_load_s.chunkat(warp_m, iv_warp_k);
          mb = mma.load rhs_load_s.chunkat(warp_n, iv_warp_k);
          mma.row.row mc, ma, mb;
        }
      }
      mma.store mc, output_s.subspan(MMA_M, MMA_N).at(warp_m, warp_n);
    }
    dma.copy output_s => output.subspan(TILE_M, TILE_N).at(block_m, block_n);
  }
}
\end{lstlisting}

The function returns \co@void@ because it writes into the \co@output@ parameter. The caller supplies that buffer and reads it after the operation completes. Inside the function, \co@output_s@ is the staging buffer used between \co@mma.store@ and the final copy.

The ceiling division \co@cdiv(M, TILE_M)@ counts the tiles along M, including a possible partial tile. The pair \co@(block_m, block_n)@ identifies the output tile assigned to the current block. The operand selections use the same output coordinates with the current K-slice coordinate.

When the matrix dimensions and the tile dimensions are all 16, the block grid is 1 by 1. Larger matrices produce more output tiles. The \co@MMA_M@, \co@MMA_N@, and \co@MMA_K@ names in the listing must be defined consistently with the intended fragment geometry before compiling it.

Section 3.5 explains how \co@subspan@ and \co@chunkat@ select tiles.

\begin{keypoint}[Quick `subspan` syntax reminder]
\texttt{subspan(h, w)} declares a fixed-size rectangular view shape, and \texttt{.at(i, j)} selects which view instance you want from the parent tensor. In other words, \texttt{lhs.\allowbreak{}subspan(TILE\_\allowbreak{}M, TILE\_\allowbreak{}K).\allowbreak{}at(block\_\allowbreak{}m, iv\_\allowbreak{}k)} means "take the \texttt{[TILE\_\allowbreak{}M, TILE\_\allowbreak{}K]} tile at \texttt{(block\_\allowbreak{}m, iv\_\allowbreak{}k)}." Use \texttt{subspan} when you want an explicit tile shape; use \texttt{chunkat} when you want to slice by equal chunks of an existing dimension.

\end{keypoint}

The accumulator spans the entire K loop. Initializing it inside that loop would discard earlier contributions, while storing inside the loop would expose only a partial result. Keeping initialization before the loop and the store after it makes the output tile's lifetime visible.
